\documentclass[conference]{IEEEtran}
\IEEEoverridecommandlockouts
\usepackage{cite}
\usepackage{amsmath,amssymb,amsfonts}
\usepackage{algorithmic}
\usepackage{graphicx}
\usepackage{textcomp}
\usepackage{xcolor}
\def\BibTeX{{\rm B\kern-.05em{\sc i\kern-.025em b}\kern-.08em
    T\kern-.1667em\lower.7ex\hbox{E}\kern-.125emX}}
\begin{document}

\title{Adaptive Secure Communication System for Unreliable and Adversarial Networks}

\author{
\IEEEauthorblockN{1\textsuperscript{st} Given Name Surname}
\IEEEauthorblockA{\textit{Department Name} \\
\textit{Institution Name}\\
City, Country \\
email@institution.edu}
\and
\IEEEauthorblockN{2\textsuperscript{nd} Given Name Surname}
\IEEEauthorblockA{\textit{Department Name} \\
\textit{Institution Name}\\
City, Country \\
email@institution.edu}
\and
\IEEEauthorblockN{3\textsuperscript{rd} Given Name Surname}
\IEEEauthorblockA{\textit{Department Name} \\
\textit{Institution Name}\\
City, Country \\
email@institution.edu}
}

\maketitle

\begin{abstract}
Secure messaging over unreliable and adversarial networks is undermined by static security
configurations that cannot respond to changing threat conditions. We present an adaptive
end-to-end encrypted chat system that combines TLS 1.3 transport security with a
Signal-protocol Double Ratchet at the application layer, dual-path TCP and UDP delivery,
and a three-state adaptive engine that automatically reconfigures cryptographic and
transport parameters in response to live network metrics and detected security events.
A lightweight per-IP intrusion detection system feeds authentication failures and replay
detections directly into the engine state machine, triggering escalation from Normal to
Unstable or High-Risk mode with immediate effect. In High-Risk mode, all payloads are
padded to a constant 4096 bytes, DH ratchet steps occur on every message, and randomized
retry delays frustrate traffic analysis. The system provides forward secrecy and
break-in recovery per the Double Ratchet specification, multipath deduplication via a
128-bit message-ID ring buffer, and persistent offline message delivery. Evaluation
confirms correct state transitions under injected packet loss and authentication failure
scenarios, with all 13 functional tests passing.
\end{abstract}

\begin{IEEEkeywords}
end-to-end encryption, adaptive security, Double Ratchet, TLS 1.3, intrusion detection,
multipath transport, forward secrecy, secure messaging
\end{IEEEkeywords}

% ============================================================
%  SECTION I — INTRODUCTION  (Member 1)
% ============================================================
\section{Introduction}

Secure communication over unreliable and adversarial networks remains an open challenge.
Standard TLS protects the channel but cannot react to rising packet loss, replay attacks,
or authentication brute-force attempts. This paper presents a multi-client end-to-end
encrypted chat system written in C (C11, POSIX) that addresses this gap through
an adaptive three-state security engine tightly coupled to both network quality metrics
and a lightweight intrusion detection system.

The main contributions of this work are:
\begin{enumerate}
    \item A three-state adaptive engine (Normal / Unstable / High-Risk) driven by live
          packet loss, round-trip time, authentication failure count, and replay detection
          count.
    \item Signal Double Ratchet at the application layer, providing per-message forward
          secrecy and break-in recovery independent of TLS.
    \item Dual-path (TCP+UDP) multipath delivery with 128-bit message-ID deduplication.
    \item A lightweight per-IP IDS that feeds security events directly into the adaptive
          state machine.
    \item Persistent offline message queue with chronological replay on reconnect.
\end{enumerate}

% ============================================================
%  SECTION II — RELATED WORK  (Member 1)
% ============================================================
\section{Related Work}

The Signal Double Ratchet Algorithm~\cite{b1} provides per-message forward secrecy and
break-in recovery through a combination of symmetric-key ratcheting and Diffie--Hellman
ratcheting. Formal security analysis of the Signal protocol was conducted by
Cohn-Gordon et al.~\cite{b2}. TLS 1.3~\cite{b3} secures the transport channel with
authenticated encryption and perfect forward secrecy at the session level. HKDF~\cite{b4}
provides the key derivation foundation used throughout our cryptographic stack.
Group messaging security is addressed by the Messaging Layer Security (MLS)
protocol~\cite{b5}, which we identify as a direction for future work.

% ============================================================
%  SECTION III — SYSTEM ARCHITECTURE  (Member 1)
% ============================================================
\section{System Architecture}

The system follows a two-tier client-server architecture. The server is a single POSIX
process that spawns a detached \texttt{pthread} per accepted client connection. A shared
client table (\texttt{ClientEntry[50]}) is protected by a \texttt{pthread\_mutex\_t}.
Two background threads run independently: an adaptive engine thread (1-second evaluation
tick) and a UDP receiver thread.

Each client runs two threads: a send thread that drains a priority queue, encrypts
messages, and transmits over both TCP and UDP; and a receive thread that deduplicates,
decrypts, and delivers messages to the UI layer.

% ============================================================
%  SECTION IV — WIRE PROTOCOL AND CONNECTION LIFECYCLE  (Member 2)
% ============================================================
\section{Wire Protocol and Connection Lifecycle}

Every message uses a uniform 28-byte binary header:

\begin{equation}
\underbrace{\texttt{ver}(1)}_{}\ |\ \underbrace{\texttt{type}(1)}_{}\ |\ \underbrace{\texttt{pri}(1)}_{}\ |\ \underbrace{\texttt{flags}(1)}_{}\ |\ \underbrace{\texttt{msg\_id}(16)}_{}\ |\ \underbrace{\texttt{len}(4)}_{}\ |\ \underbrace{\texttt{crc32}(4)}_{}
\end{equation}

The 128-bit \texttt{msg\_id} is used for deduplication across TCP and UDP paths.
All \texttt{MSG\_CHAT} payloads are padded to a constant 4112 bytes
(16-byte IV + 4096-byte AES ciphertext) to resist traffic analysis.

% ============================================================
%  SECTION V — CRYPTOGRAPHIC DESIGN  (Member 2)
% ============================================================
\section{Cryptographic Design}

\subsection{Cryptographic Stack}

The system employs a four-layer cryptographic stack:

\begin{enumerate}
    \item \textbf{Key Agreement} --- X25519 ECDH (\texttt{MSG\_DH\_INIT} /
          \texttt{MSG\_DH\_RESP}). The shared secret is expanded via HKDF~\cite{b4}
          to produce initial ratchet keys.
    \item \textbf{Identity Authentication} --- RSA-2048 signatures over the username
          bytes. The server stores the public key on first login and verifies it on
          subsequent logins.
    \item \textbf{Transport Confidentiality} --- TLS 1.3~\cite{b3} with TLS 1.2
          and below disabled.
    \item \textbf{E2EE Confidentiality} --- AES-256-CBC with per-message keys derived
          from the Double Ratchet chain. Operates inside TLS so that TLS compromise
          does not expose message content.
\end{enumerate}

\subsection{Double Ratchet Protocol}

The Double Ratchet~\cite{b1} maintains two interlocked ratchets. The symmetric ratchet
advances on every message:
\begin{align}
    \textit{msg\_key}       &= \text{HMAC-SHA256}(\textit{chain\_key},\ \texttt{0x01}) \\
    \textit{new\_chain\_key} &= \text{HMAC-SHA256}(\textit{chain\_key},\ \texttt{0x02})
\end{align}
Each message uses a unique derived key; the old chain key is discarded immediately.
The DH ratchet injects fresh Diffie--Hellman entropy every \texttt{dh\_ratchet\_freq}
messages, providing break-in recovery. In High-Risk mode this frequency is set to 1,
meaning a new X25519 keypair is generated per message.

All key material is erased with \texttt{OPENSSL\_cleanse()} immediately after use.

% ============================================================
%  SECTION VI — ADAPTIVE ENGINE AND IDS  (Member 3)
% ============================================================
\section{Adaptive Engine and Intrusion Detection}

\subsection{Three-State Adaptive Engine}

A background thread evaluates a \texttt{Metrics} structure every second and updates the
shared engine state read by all modules. Three operating modes are defined:

\begin{itemize}
    \item \textbf{Normal} --- standard retries (3), 100\,ms retry delay, DH ratchet
          every 10 messages.
    \item \textbf{Unstable} --- increased retries (7), 200\,ms delay; triggered by
          $\geq$5\% packet loss or $\geq$3 consecutive timeouts.
    \item \textbf{High-Risk} --- maximum retries (10), 300\,ms base delay, forced
          4096-byte padding, random 100--500\,ms delays, DH ratchet every message;
          triggered immediately by $\geq$20\% packet loss, $\geq$5 auth failures, or
          $\geq$3 replay detections.
\end{itemize}

Escalation is immediate. Downgrade to Normal requires 30 seconds of stable metrics
(hysteresis to prevent state thrashing).

\subsection{Intrusion Detection System}

The IDS maintains a table of 256 per-IP records. After five authentication failures
from the same IP, the address is blocked for 300 seconds and
\texttt{metrics\_record\_auth\_fail()} is called, escalating the engine to High-Risk.
Replay detections likewise call \texttt{metrics\_record\_replay()}. Blocked IPs are
rejected before the TLS handshake completes.

% ============================================================
%  SECTION VII — MULTIPATH TRANSPORT AND OFFLINE QUEUE  (Member 3)
% ============================================================
\section{Multipath Transport and Offline Message Queue}

\subsection{Dual-Path Delivery}

Every outgoing message is transmitted over both the TLS/TCP primary path and a UDP
backup path simultaneously. A 1024-slot deduplication ring buffer keyed on the 128-bit
\texttt{msg\_id} discards duplicate deliveries silently. The ring buffer is
\texttt{O(1024)} at typical chat rates and protected by a \texttt{pthread\_mutex\_t}.

\subsection{Priority Queue}

The client send thread drains a three-level priority queue: Normal, Urgent, and
Critical. Critical messages wake the send thread immediately and receive two additional
multipath retry attempts.

\subsection{Offline Message Queue}

Messages addressed to a disconnected recipient are stored on disk under
\texttt{data/offline\_queue/<username>/} as one file per message, named by millisecond
timestamp and message ID for lexicographic chronological ordering. Plaintext is stored
(not ciphertext) because the recipient's future ratchet key is unknown at queue time;
re-encryption occurs when the recipient reconnects. Queue capacity is 500 messages per
user. Directories are mode \texttt{0700} and files \texttt{0600}; usernames are
validated against \texttt{[a-zA-Z0-9\_-]} to prevent directory traversal.

% ============================================================
%  SECTION VIII — EVALUATION  (Member 3)
% ============================================================
\section{Evaluation}

\subsection{Functional Correctness}

All 13 functional tests pass. Key results are summarised below:

\begin{itemize}
    \item \textbf{Ratchet key uniqueness}: 100 derived message keys are all distinct.
    \item \textbf{Forward secrecy}: deleting ratchet state at message 50 leaves all
          prior messages unrecoverable.
    \item \textbf{Break-in recovery}: exposing the chain key at message 50; messages
          after the subsequent DH ratchet step are unreadable to an attacker holding
          that chain key.
    \item \textbf{Constant payload size}: all \texttt{MSG\_CHAT} frames are 4112 bytes
          as confirmed via \texttt{tcpdump}.
    \item \textbf{Multipath dedup}: the same \texttt{msg\_id} arriving via TCP and UDP
          is processed exactly once.
    \item \textbf{Engine transitions}: injecting 6\% packet loss triggers
          Normal$\to$Unstable; 21\% triggers Unstable$\to$High-Risk; 5 auth failures
          trigger immediate escalation to High-Risk from any state.
    \item \textbf{Hysteresis}: 30 seconds of stable metrics after High-Risk produces
          a return to Normal.
    \item \textbf{IDS block}: 5 auth failures from one IP blocks it; TLS handshake
          refused. Block lifted after 300 seconds.
\end{itemize}

% ============================================================
%  SECTION IX — LIMITATIONS AND FUTURE WORK  (Member 2)
% ============================================================
\section{Limitations and Future Work}

The current design has several acknowledged limitations. First, the server acts as a
trusted re-encryption relay: each \texttt{MSG\_CHAT} is decrypted with the sender's
ratchet and re-encrypted with the recipient's ratchet, so plaintext passes through
server memory. Pure end-to-end encryption in the Signal sense would require X3DH
pre-key exchange~\cite{b6} so the server never sees plaintext. Second, offline queued
messages are stored as plaintext, exposing them in the event of a server breach.
Third, TLS certificates are self-signed, suitable for lab evaluation but not
production deployment. Fourth, RSA keypairs are stored as plain files rather than
in hardware security modules. Fifth, broadcast messaging is implemented as a loop of
per-client directed encryptions rather than a proper group ratchet such as
MLS~\cite{b5}. Addressing these limitations, particularly adopting X3DH and MLS,
constitutes planned future work.

% ============================================================
%  SECTION X — CONCLUSION  (Member 1)
% ============================================================
\section{Conclusion}

This paper has presented an adaptive secure communication system that couples live
network quality and security event metrics to cryptographic and transport configuration
through a three-state engine. The system demonstrates that security posture need not
be static: escalation to High-Risk mode automatically enables per-message DH ratchet
steps, constant-size padding, and randomized retry delays in response to detected
threats or network degradation. All functional correctness tests pass, including
verification of forward secrecy, break-in recovery, constant payload size, and correct
engine state transitions. Future work will address the server trust limitation via
X3DH pre-key exchange and extend group messaging to a proper group ratchet.

\section*{Acknowledgment}

\begin{thebibliography}{00}
\bibitem{b1} M. Marlinspike and T. Perrin, ``The Double Ratchet Algorithm,''
    Signal Foundation, 2016. [Online]. Available: https://signal.org/docs/specifications/doubleratchet/

\bibitem{b2} K. Cohn-Gordon, C. Cremers, B. Dowling, L. Garratt, and D. Stebila,
    ``A Formal Security Analysis of the Signal Messaging Protocol,''
    in \textit{Proc. IEEE European Symposium on Security and Privacy (EuroS\&P)},
    Paris, France, 2017, pp. 451--466.

\bibitem{b3} E. Rescorla, ``The Transport Layer Security (TLS) Protocol Version 1.3,''
    IETF RFC 8446, Aug. 2018.

\bibitem{b4} H. Krawczyk and P. Eronen, ``HMAC-based Extract-and-Expand Key
    Derivation Function (HKDF),'' IETF RFC 5869, May 2010.

\bibitem{b5} R. Barnes, B. Beurdouche, R. Robert, J. Millican, E. Omara, and K. Cohn-Gordon,
    ``The Messaging Layer Security (MLS) Protocol,''
    IETF RFC 9420, Jul. 2023.

\bibitem{b6} M. Marlinspike and T. Perrin, ``The X3DH Key Agreement Protocol,''
    Signal Foundation, 2016. [Online]. Available: https://signal.org/docs/specifications/x3dh/
\end{thebibliography}

\end{document}
